% ============================================================
\chapter{Processus Stochastiques --- D\'efinitions}
\label{ch:definitions}
% ============================================================

Qu'est-ce qu'un processus stochastique~? Intuitivement, c'est une quantit\'e
al\'eatoire qui \'evolue dans le temps~: le cours d'une action, la position
d'une particule, le nombre de clients dans une file d'attente. Mais
formellement, c'est un objet math\'ematique bien plus riche qu'une simple
variable al\'eatoire~: c'est une \emph{famille} de variables al\'eatoires
index\'ees par le temps, et l'on peut l'\'etudier soit \og \`a temps fix\'e\fg{}
(les lois fini-dimensionnelles), soit \og \`a \'ev\'enement fix\'e\fg{} (les
trajectoires). Ce double point de vue est au c\oe ur de toute la th\'eorie.

% ------------------------------------------------------------
\section{D\'efinition d'un processus stochastique}
% ------------------------------------------------------------

\begin{definition}[Processus stochastique]
Soit $(\Omega, \mathcal{F}, \PP)$ un espace de probabilit\'e et $T$ un ensemble
d'indices (typiquement $T = \N$, $\Z_+$ ou $\R_+$). Un \emph{processus stochastique}
\`a valeurs dans un espace mesurable $(E, \mathcal{E})$ est une famille de v.a.\
\[
  X = (X_t)_{t \in T}, \quad X_t : \Omega \to E,
\]
o\`u chaque $X_t$ est $\mathcal{F}/\mathcal{E}$-mesurable.

De mani\`ere \'equivalente, $X$ est une application mesurable
$X : (T \times \Omega, \mathcal{B}(T) \otimes \mathcal{F}) \to (E, \mathcal{E})$
(lorsque la mesurabilit\'e jointe est requise).
\end{definition}

\begin{definition}[Trajectoire]
Pour $\omega \in \Omega$ fix\'e, l'application $t \mapsto X_t(\omega)$ est appel\'ee
la \emph{trajectoire} (ou \emph{r\'ealisation}, ou \emph{chemin d'\'echantillon})
du processus associ\'ee \`a $\omega$.
\end{definition}

\begin{exemple}
\begin{enumerate}
  \item \textbf{Marche al\'eatoire simple}~: soit $(Y_n)_{n \geq 1}$ une suite i.i.d.\ avec
    $\PP(Y_n = 1) = \PP(Y_n = -1) = 1/2$. Le processus $S_n = \sum_{k=1}^{n} Y_k$
    pour $n \geq 0$ (avec $S_0 = 0$) est un processus stochastique \`a temps discret
    et \`a valeurs dans $\Z$.
  \item \textbf{Processus de Poisson}~: $N = (N_t)_{t \geq 0}$ \`a valeurs dans $\N$,
    \`a trajectoires croissantes et c\`adl\`ag.
  \item \textbf{Mouvement brownien}~: $B = (B_t)_{t \geq 0}$ \`a valeurs dans $\R$,
    \`a trajectoires continues.
\end{enumerate}
\end{exemple}

\section{Classification des processus}

\begin{definition}[Temps discret vs.\ continu]
\begin{itemize}
  \item Un processus est \`a \emph{temps discret} si $T$ est d\'enombrable
    (par exemple $T = \N$ ou $T = \Z$).
  \item Un processus est \`a \emph{temps continu} si $T$ est un intervalle de $\R$
    (par exemple $T = \R_+$ ou $T = [0, T_0]$).
\end{itemize}
\end{definition}

\begin{definition}[Espace d'\'etats]
L'espace $(E, \mathcal{E})$ est appel\'e l'\emph{espace d'\'etats} du processus.
\begin{itemize}
  \item Si $E$ est d\'enombrable (fini ou infini), on parle de processus \`a
    \emph{espace d'\'etats discret}.
  \item Si $E = \R^d$, on parle de processus \`a \emph{espace d'\'etats continu}.
\end{itemize}
\end{definition}

\begin{definition}[R\'egularit\'e des trajectoires]
Un processus $(X_t)_{t \geq 0}$ \`a valeurs r\'eelles est dit~:
\begin{itemize}
  \item \emph{continu} si p.s.\ $t \mapsto X_t(\omega)$ est continue,
  \item \emph{c\`adl\`ag} (continu \`a droite avec limites \`a gauche) si p.s.\
    pour tout $t$, $X_t = X_{t+}$ et $X_{t-}$ existe,
  \item \emph{c\`agl\`ad} (continu \`a gauche avec limites \`a droite) si p.s.\
    $X_t = X_{t-}$ et $X_{t+}$ existe.
\end{itemize}
\end{definition}

\section{Filtrations}

\begin{definition}[Filtration]
\label{def:filtration}
Une \emph{filtration} sur $(\Omega, \mathcal{F}, \PP)$ est une famille croissante
de sous-tribus $\mathbb{F} = (\mathcal{F}_t)_{t \in T}$~:
\[
  s \leq t \implies \mathcal{F}_s \subset \mathcal{F}_t \subset \mathcal{F}.
\]
L'espace $(\Omega, \mathcal{F}, \mathbb{F}, \PP)$ est appel\'e \emph{espace de
probabilit\'e filtr\'e}.
\end{definition}

\begin{intuition}[title=\textbf{Information croissante}]
La filtration mod\'elise l'\emph{information disponible au cours du temps}.
La tribu $\mathcal{F}_t$ repr\'esente l'ensemble des \'ev\'enements que l'on peut
observer (d\'ecider s'ils sont r\'ealis\'es ou non) \`a l'instant $t$. La croissance
de la filtration traduit le fait que l'information ne se perd pas~: ce qu'on sait
\`a l'instant $s$ reste connu \`a l'instant $t > s$.
\end{intuition}

\begin{definition}[Filtration naturelle]
La \emph{filtration naturelle} (ou \emph{filtration engendr\'ee}) par un processus
$(X_t)_{t \in T}$ est
\[
  \mathcal{F}_t^X = \sigma(X_s : s \leq t), \quad t \in T.
\]
C'est la plus petite filtration par rapport \`a laquelle $(X_t)$ est adapt\'e.
\end{definition}

\begin{definition}[Conditions habituelles]
\label{def:conditions_habituelles}
Une filtration $(\mathcal{F}_t)_{t \geq 0}$ satisfait les \emph{conditions
habituelles} (de Dellacherie--Meyer) si~:
\begin{enumerate}
  \item \textbf{Compl\'etude}~: $\mathcal{F}_0$ contient tous les ensembles
    $\PP$-n\'egligeables de $\mathcal{F}$.
  \item \textbf{Continuit\'e \`a droite}~: pour tout $t \geq 0$,
    $\mathcal{F}_t = \mathcal{F}_{t+} := \bigcap_{s > t} \mathcal{F}_s$.
\end{enumerate}
\end{definition}

\begin{remarque}
On peut toujours \og compl\'eter \fg et \og r\'egulariser \fg une filtration.
Soit $\mathcal{N}$ la collection des ensembles $\PP$-n\'egligeables. On d\'efinit
\[
  \overline{\mathcal{F}}_t = \sigma(\mathcal{F}_t \cup \mathcal{N}), \quad
  \widehat{\mathcal{F}}_t = \bigcap_{s > t} \overline{\mathcal{F}}_s.
\]
La filtration $(\widehat{\mathcal{F}}_t)$ satisfait les conditions habituelles.
\end{remarque}

\section{Processus adapt\'es et pr\'evisibles}

\begin{definition}[Processus adapt\'e]
Un processus $(X_t)_{t \in T}$ est dit \emph{adapt\'e} \`a la filtration
$(\mathcal{F}_t)_{t \in T}$ si pour tout $t \in T$, la variable $X_t$ est
$\mathcal{F}_t$-mesurable.
\end{definition}

\begin{definition}[Processus pr\'evisible --- temps discret]
En temps discret, un processus $(X_n)_{n \geq 0}$ est \emph{pr\'evisible}
par rapport \`a $(\mathcal{F}_n)$ si $X_0$ est d\'eterministe (ou $\mathcal{F}_0$-mesurable)
et $X_n$ est $\mathcal{F}_{n-1}$-mesurable pour tout $n \geq 1$.
\end{definition}

\begin{definition}[Tribu pr\'evisible --- temps continu]
\label{def:tribu_previsible}
En temps continu, la \emph{tribu pr\'evisible} $\mathcal{P}$ sur
$\R_+ \times \Omega$ est engendr\'ee par les ensembles de la forme
$(s, t] \times F$ avec $0 \leq s < t$ et $F \in \mathcal{F}_s$, et les ensembles
$\{0\} \times F_0$ avec $F_0 \in \mathcal{F}_0$.

Un processus est \emph{pr\'evisible} s'il est mesurable par rapport \`a $\mathcal{P}$.
\end{definition}

\begin{remarque}
Tout processus adapt\'e continu \`a gauche est pr\'evisible.
Tout processus pr\'evisible est adapt\'e.
En revanche, un processus adapt\'e c\`adl\`ag n'est pas n\'ecessairement pr\'evisible.
\end{remarque}

\section{Temps d'arr\^et}

\begin{definition}[Temps d'arr\^et]
\label{def:temps_arret}
Soit $(\mathcal{F}_t)_{t \in T}$ une filtration. Une v.a.\ $\tau : \Omega \to
T \cup \{+\infty\}$ est un \emph{temps d'arr\^et} (ou \emph{temps de Markov}) si~:
\begin{itemize}
  \item En temps discret ($T = \N$)~: $\{\tau = n\} \in \mathcal{F}_n$ pour tout $n \in \N$.
  \item En temps continu ($T = \R_+$)~: $\{\tau \leq t\} \in \mathcal{F}_t$ pour tout $t \geq 0$.
\end{itemize}
\end{definition}

\begin{remarque}
En temps discret, la condition $\{\tau = n\} \in \mathcal{F}_n$ est \'equivalente
\`a $\{\tau \leq n\} \in \mathcal{F}_n$. En temps continu, si la filtration est
continue \`a droite, la condition $\{\tau < t\} \in \mathcal{F}_t$ est \'equivalente.
\end{remarque}

\begin{exemple}
\begin{enumerate}
  \item Toute constante $\tau = t_0$ est un temps d'arr\^et.
  \item Le \emph{temps de premier passage} $\tau_a = \inf\{t \geq 0 : X_t = a\}$
    est un temps d'arr\^et si le processus $(X_t)$ est adapt\'e et a des trajectoires
    continues (ou c\`adl\`ag sous les conditions habituelles).
  \item Le \emph{temps de premier entr\'ee} dans un ouvert $G$~:
    $\tau_G = \inf\{t \geq 0 : X_t \in G\}$ est un temps d'arr\^et sous les
    conditions habituelles.
\end{enumerate}
\end{exemple}

\begin{definition}[Tribu des \'ev\'enements ant\'erieurs \`a $\tau$]
\label{def:tribu_tau}
Si $\tau$ est un temps d'arr\^et, la \emph{tribu des \'ev\'enements ant\'erieurs \`a
$\tau$} est
\[
  \mathcal{F}_{\tau} = \{A \in \mathcal{F} : \forall t \in T,\; A \cap \{\tau \leq t\}
  \in \mathcal{F}_t\}.
\]
\end{definition}

\begin{proposition}[Propri\'et\'es des temps d'arr\^et]
Soient $\sigma, \tau$ des temps d'arr\^et.
\begin{enumerate}
  \item $\sigma \wedge \tau = \min(\sigma, \tau)$ et $\sigma \vee \tau = \max(\sigma, \tau)$
    sont des temps d'arr\^et.
  \item Si $\sigma \leq \tau$ p.s., alors $\mathcal{F}_{\sigma} \subset \mathcal{F}_{\tau}$.
  \item En temps discret, $\tau + 1$ est un temps d'arr\^et.
  \item $\{\sigma \leq \tau\}$, $\{\sigma < \tau\}$, $\{\sigma = \tau\}$ appartiennent
    \`a $\mathcal{F}_{\sigma} \cap \mathcal{F}_{\tau}$.
\end{enumerate}
\end{proposition}

\begin{definition}[Processus arr\^et\'e]
Le \emph{processus arr\^et\'e} au temps d'arr\^et $\tau$ est
\[
  X_t^{\tau} = X_{t \wedge \tau} = X_{\min(t, \tau)}, \quad t \in T.
\]
\end{definition}

\section{Lois de dimension finie}

\begin{definition}[Lois de dimension finie]
Les \emph{lois de dimension finie} (ou \emph{distributions fini-dimensionnelles})
d'un processus $(X_t)_{t \in T}$ sont les lois des vecteurs al\'eatoires
$(X_{t_1}, \ldots, X_{t_n})$ pour tout choix $t_1 < \cdots < t_n$ dans $T$ et
tout $n \geq 1$.
\end{definition}

\begin{theoreme}[Th\'eor\`eme d'extension de Kolmogorov]
\label{thm:kolmogorov_extension}
Soit $(E, \mathcal{E})$ un espace polonais. Soit $(\mu_{t_1, \ldots, t_n})$
une famille de mesures de probabilit\'e sur $(E^n, \mathcal{E}^{\otimes n})$
satisfaisant les conditions de coh\'erence~:
\begin{enumerate}
  \item \textbf{Invariance par permutation}~: pour toute permutation $\pi$ de
    $\{1, \ldots, n\}$,
    \[
      \mu_{t_{\pi(1)}, \ldots, t_{\pi(n)}}(B_{\pi(1)} \times \cdots \times B_{\pi(n)})
      = \mu_{t_1, \ldots, t_n}(B_1 \times \cdots \times B_n).
    \]
  \item \textbf{Marginalisation}~:
    $\mu_{t_1, \ldots, t_n}(B_1 \times \cdots \times B_n \times E) = \mu_{t_1, \ldots, t_n}(B_1 \times \cdots \times B_n)$.
\end{enumerate}
Alors il existe un espace de probabilit\'e $(\Omega, \mathcal{F}, \PP)$ et un
processus $(X_t)_{t \in T}$ sur cet espace dont les lois fini-dimensionnelles
sont les $(\mu_{t_1, \ldots, t_n})$.
\end{theoreme}

\begin{attention}[title=\textbf{Limites du th\'eor\`eme de Kolmogorov}]
Le th\'eor\`eme de Kolmogorov garantit l'existence d'un processus avec les bonnes
lois fini-dimensionnelles, mais \emph{pas} la r\'egularit\'e des trajectoires.
Des propri\'et\'es comme la continuit\'e des trajectoires n\'ecessitent des r\'esultats
suppl\'ementaires (par exemple le th\'eor\`eme de Kolmogorov--\v{C}entsov).
\end{attention}

\section{Th\'eor\`eme de Kolmogorov--\v{C}entsov}

\begin{theoreme}[Kolmogorov--\v{C}entsov]
\label{thm:kolmogorov_centsov}
Soit $(X_t)_{t \in [0,T]}$ un processus stochastique \`a valeurs r\'eelles.
S'il existe des constantes $\alpha > 0$, $\beta > 0$ et $C > 0$ telles que
\[
  \E[\abs{X_t - X_s}^{\alpha}] \leq C \abs{t - s}^{1 + \beta}, \quad \forall s, t \in [0, T],
\]
alors il existe une modification $(\tilde{X}_t)$ de $(X_t)$ dont les trajectoires
sont p.s.\ h\"old\'eriennes d'exposant $\gamma$ pour tout $\gamma < \beta/\alpha$.
\end{theoreme}

\begin{remarque}
Rappelons qu'une \emph{modification} de $(X_t)$ est un processus $(\tilde{X}_t)$
d\'efini sur le m\^eme espace de probabilit\'e tel que $\PP(X_t = \tilde{X}_t) = 1$
pour tout $t$.
\end{remarque}

\section{Processus gaussiens}

\begin{definition}[Processus gaussien]
Un processus $(X_t)_{t \in T}$ est \emph{gaussien} si pour tout choix
$t_1, \ldots, t_n \in T$, le vecteur $(X_{t_1}, \ldots, X_{t_n})$ suit une loi
gaussienne multivari\'ee.

Un processus gaussien est enti\`erement d\'etermin\'e par sa \emph{fonction moyenne}
$m(t) = \E[X_t]$ et sa \emph{fonction de covariance}
$K(s, t) = \Cov(X_s, X_t)$.
\end{definition}

\begin{exemple}
Le mouvement brownien standard est le processus gaussien avec $m(t) = 0$ et
$K(s, t) = \min(s, t)$. Le \emph{pont brownien} sur $[0, 1]$ est le processus
gaussien avec $m(t) = 0$ et $K(s, t) = \min(s, t) - st$.
\end{exemple}

\section{Stationnarit\'e et ergodicit\'e}

\begin{definition}[Stationnarit\'e forte]
Un processus $(X_t)_{t \in T}$ est \emph{strictement stationnaire} si pour tout
$h > 0$ et tout $n \geq 1$,
\[
  (X_{t_1+h}, \ldots, X_{t_n+h}) \overset{\mathcal{L}}{=} (X_{t_1}, \ldots, X_{t_n}).
\]
\end{definition}

\begin{definition}[Stationnarit\'e faible]
Un processus $(X_t)$ de carr\'e int\'egrable est \emph{faiblement stationnaire}
(ou \emph{stationnaire au second ordre}, ou \emph{stationnaire au sens large}) si~:
\begin{enumerate}
  \item $\E[X_t] = \mu$ est constante,
  \item $\Cov(X_s, X_t) = R(t - s)$ ne d\'epend que de $t - s$.
\end{enumerate}
La fonction $R$ est appel\'ee la \emph{fonction d'autocorr\'elation}.
\end{definition}

\section{Exercices}

\begin{exercice}
Montrer que la filtration naturelle d'un processus continu satisfait la
condition de continuit\'e \`a droite apr\`es compl\'etion.
\end{exercice}

\begin{exercice}
Soit $\tau$ un temps d'arr\^et par rapport \`a $(\mathcal{F}_n)_{n \geq 0}$.
Montrer que $X_{\tau}$ est $\mathcal{F}_{\tau}$-mesurable lorsque le processus
$(X_n)$ est adapt\'e.
\end{exercice}

\begin{exercice}
Soit $(B_t)_{t \geq 0}$ un mouvement brownien standard et $a > 0$. Montrer que
$\tau_a = \inf\{t \geq 0 : B_t = a\}$ est un temps d'arr\^et par rapport \`a la
filtration naturelle compl\'et\'ee de $B$.
\end{exercice}

\begin{exercice}
V\'erifier que le pont brownien $(X_t)_{t \in [0,1]}$ d\'efini par
$X_t = B_t - tB_1$ est un processus gaussien et calculer sa fonction de covariance.
\end{exercice}

\begin{exercice}
Soit $(X_t)_{t \geq 0}$ un processus v\'erifiant les hypoth\`eses du th\'eor\`eme de
Kolmogorov--\v{C}entsov avec $\alpha = 4$ et $\beta = 1$. Quels sont les exposants
de H\"older possibles pour la modification continue ?
\end{exercice}
